% Guide (ohne Nummer) – From Topic to Thesis: A W-Seminar Research Guide
% Leitfaden im Handbuch-Layout (Abschnitte statt Boxen), nach einer Vorlage der Lehrkraft.
% Übersetzen: latexmk -xelatex guide-question-to-thesis.tex
% Begriffe wie W4–W7 und research-question.html (spark question → research question → working thesis
% → testable claim → thesis statement, sechs Marker); FINER als zweite Sicht auf dieselben Kriterien.
% ============================================================
% ab-vorlage.tex – gemeinsame Vorlage der Arbeitsblätter
% englisch/w-seminar-time/arbeitsblaetter/  (W-Seminar "Time")
% Übersetzen mit:  latexmk -xelatex <datei>.tex   (oder zweimal xelatex)
% Gestaltung wie englisch/13/questions-on-the-text/arbeitsblaetter/ab-vorlage.tex,
% Antwortplatz-Makros (8-mm-Linien) wie informatik/10/java/arbeitsblaetter/ab-vorlage.tex
% Schriften: Carlito (Fließtext), Poppins (Überschriften)
% ============================================================
\documentclass[10pt,a4paper]{article}
\usepackage[a4paper,top=12mm,bottom=13mm,left=14mm,right=14mm,headsep=4mm,headheight=12pt,footskip=7mm]{geometry}
\usepackage{fontspec}
\setmainfont{Carlito}
\newfontfamily\headfont{Poppins}[UprightFont=* Medium, BoldFont=* Medium]
\newfontfamily\headlight{Poppins Light}
\usepackage{xcolor}
\definecolor{prim}{HTML}{2563EB}
\definecolor{primdark}{HTML}{1E3A8A}
\definecolor{sec}{HTML}{7C3AED}
\definecolor{dusk}{HTML}{171A3F}
\definecolor{brass}{HTML}{D9AB2F}
\definecolor{primtint}{HTML}{EFF6FF}
\definecolor{sectint}{HTML}{F5F3FF}
\definecolor{ink}{HTML}{1E293B}
\definecolor{muted}{HTML}{64748B}
\definecolor{rule}{HTML}{CBD5E1}
\definecolor{good}{HTML}{059669}
\definecolor{goodtint}{HTML}{ECFDF5}
\definecolor{warn}{HTML}{B45309}
\definecolor{warntint}{HTML}{FFFBEB}
\definecolor{bad}{HTML}{B91C1C}
\definecolor{badtint}{HTML}{FEF2F2}
% Farben der Zeit-Dimensionen (wie auf der Website)
\definecolor{dlit}{HTML}{B83280}
\definecolor{dphys}{HTML}{2563EB}
\definecolor{dcs}{HTML}{0F8F84}
\definecolor{dpsy}{HTML}{C26A06}
\definecolor{dlang}{HTML}{7C3AED}
\definecolor{dbey}{HTML}{4B5563}
\color{ink}
\usepackage[most]{tcolorbox}
\usepackage{tikz}
\usetikzlibrary{positioning,calc,arrows.meta,shapes.misc}
\usepackage{enumitem}
\setlist{nosep,leftmargin=1.1em,topsep=1pt}
\setlist[itemize]{label=\textcolor{prim}{\small\textbullet}}
\usepackage{tabularx,array,booktabs,colortbl}
\usepackage{multicol}
\usepackage{fancyhdr}
\usepackage{lastpage}
\usepackage{amssymb}
\usepackage{calc}
\usepackage{qrcode}
\usepackage{microtype}
\setlength{\parindent}{0pt}
\setlength{\parskip}{0pt}
\linespread{1.02}

% ---------- Kopf- und Fußzeile ----------
\newcommand{\wskopf}{W-Seminar English · Time · Year 12}
\pagestyle{fancy}
\fancyhf{}
\renewcommand{\headrulewidth}{0pt}
\fancyhead[L]{\footnotesize\headlight\textcolor{muted}{\wskopf}}
\fancyhead[R]{\ifnum\value{page}=1\footnotesize\headlight\textcolor{muted}{Name: \rule{38mm}{0.3pt}\quad Date: \rule{18mm}{0.3pt}}\fi}
\fancyfoot[L]{\scriptsize\textcolor{muted}{edu-mrh.de/englisch/w-seminar-time}}
\fancyfoot[R]{\scriptsize\textcolor{muted}{\thepage\,/\,\pageref*{LastPage}}}

% ---------- Titel ----------
% \wstitel{Kürzel}{Titel}{Untertitel}
\newcommand{\wstitel}[3]{%
  \begin{tcolorbox}[enhanced,colback=dusk,colframe=dusk,arc=2mm,boxrule=0pt,
     left=4mm,right=4mm,top=2.2mm,bottom=2.2mm,
     interior style={left color=dusk,right color=primdark}]
    \color{white}{\headlight\footnotesize #1}\\[0.3mm]
    {\headfont\LARGE #2}\\[0.6mm]
    {\small #3}
  \end{tcolorbox}\vspace{1mm}}

% ---------- Boxen ----------
% \begin{wsbox}[Optionen]{Nummer}{Überschrift} ... \end{wsbox}
\newtcolorbox{wsbox}[3][]{enhanced,breakable,colback=white,colframe=rule,boxrule=0.5pt,arc=1.5mm,
  left=2.5mm,right=2.5mm,top=1.3mm,bottom=1.3mm,
  title={\headfont\normalsize\textcolor{white}{\colorbox{prim}{\,#2\,}}\;\textcolor{ink}{#3}},
  coltitle=ink,colbacktitle=white,attach title to upper={\par\vspace{0.6mm}},before skip=1.4mm,after skip=1.4mm,#1}
\newtcolorbox{tintbox}[1][]{enhanced,colback=primtint,colframe=primtint,boxrule=0pt,arc=1.2mm,
  left=2mm,right=2mm,top=1.2mm,bottom=1.2mm,#1}
\newtcolorbox{secbox}[1][]{enhanced,colback=sectint,colframe=sectint,boxrule=0pt,arc=1.2mm,
  left=2mm,right=2mm,top=1.2mm,bottom=1.2mm,#1}
\newtcolorbox{goodbox}[1][]{enhanced,colback=goodtint,colframe=goodtint,boxrule=0pt,arc=1.2mm,
  left=2mm,right=2mm,top=1.2mm,bottom=1.2mm,#1}
\newtcolorbox{warnbox}[1][]{enhanced,colback=warntint,colframe=warntint,boxrule=0pt,arc=1.2mm,
  left=2mm,right=2mm,top=1.2mm,bottom=1.2mm,#1}
\newtcolorbox{badbox}[1][]{enhanced,colback=badtint,colframe=badtint,boxrule=0pt,arc=1.2mm,
  left=2mm,right=2mm,top=1.2mm,bottom=1.2mm,#1}

\newcommand{\kl}[1]{{\headfont\small\textcolor{prim}{#1}}}      % kleine Zwischenüberschrift
\newcommand{\klsec}[1]{{\headfont\small\textcolor{sec}{#1}}}
\newcommand{\klc}[2]{{\headfont\small\textcolor{#1}{#2}}}       % Zwischenüberschrift in Farbe
\newcommand{\plus}{\textcolor{good}{\textbf{+}}\ }
\newcommand{\minus}{\textcolor{bad}{\textbf{–}}\ }
\newcommand{\pfeil}{\textcolor{prim}{$\rightarrow$}\ }
\newcommand{\cb}{\raisebox{-0.2ex}{\textcolor{muted}{$\square$}}\ }   % Ankreuzkästchen

% ---------- Fragepfad: vom ersten Funken zur These (W2–W7) ----------
% Die Forschungsfrage entwickelt sich über die Blätter; jedes Blatt hat eine Box „My question – version n“.
% \fragepfad{k}: ganzer Pfad mit Blattangaben, Stufe k hervorgehoben (frühere getönt, spätere blass)
% \fragepfadklein{k}: kompakt, z. B. rechts in der Boxüberschrift (\hfill\fragepfadklein{3})
\newcommand{\pfadfarben}[2]{%
  \ifnum#1<#2\relax\def\pfc{primtint}\def\pdc{primtint}\def\ptc{prim}%
  \else\ifnum#1=#2\relax\def\pfc{prim}\def\pdc{prim}\def\ptc{white}%
  \else\def\pfc{white}\def\pdc{rule}\def\ptc{muted}\fi\fi}
\newcommand{\fragepfad}[1]{%
  \begin{tikzpicture}[x=1mm,y=1mm]
  \foreach \t/\w [count=\i] in {spark question/W2–W3,research question/{W4 · fixed by the end of 12/1},working thesis/W4–W5,testable claim/W6–W7 · for data: your hypothesis,thesis statement/the paper · backed by your results}{%
    \pfadfarben{\i}{#1}%
    \node[rounded corners=1mm,fill=\pfc,draw=\pdc,minimum width=30mm,minimum height=5.5mm,inner sep=1pt,
      text=\ptc,font=\headfont\scriptsize] (s\i) at ({(\i-1)*35},0) {\t};
    \node[font=\tiny,text=muted,below=0.3mm of s\i,inner sep=0.5pt] {\w};
    \ifnum\i>1 \draw[-{Stealth[length=1.3mm]},muted] ({(\i-1)*35-19.5},0) -- ({(\i-1)*35-15.5},0);\fi
  }
  \end{tikzpicture}}
\newcommand{\fragepfadklein}[1]{%
  \begin{tikzpicture}[baseline=-0.7ex,x=1mm,y=1mm]
  \foreach \t [count=\i] in {spark,question,working thesis,testable claim,thesis}{%
    \pfadfarben{\i}{#1}%
    \pgfmathtruncatemacro\j{\i-1}%
    \ifnum\i=1 \node[rounded corners=0.8mm,fill=\pfc,draw=\pdc,inner xsep=1mm,inner ysep=0.6mm,text=\ptc,font=\headfont\tiny] (k1) at (0,0) {\t};%
    \else \node[rounded corners=0.8mm,fill=\pfc,draw=\pdc,inner xsep=1mm,inner ysep=0.6mm,text=\ptc,font=\headfont\tiny,right=2.2mm of k\j] (k\i) {\t};%
      \draw[-{Stealth[length=1mm]},muted] (k\j) -- (k\i);\fi
  }
  \end{tikzpicture}}

% ---------- Antwortplatz (8 mm, wie informatik/10/java) ----------
\newif\ifloesung
\ifdefined\mitloesung \loesungtrue \else \loesungfalse \fi
\newcommand{\lsgfarbe}{\color{sec}\itshape}
\newlength{\zeilenabstand}\setlength{\zeilenabstand}{8mm}
\tikzset{schreiblinie/.style={draw=black!45,line width=0.35pt}}
% Lücke: \luecke[Breite] / \lsg[Breite]{Text}
\newcommand{\lsg}[2][45mm]{%
  \makebox[#1][l]{\rlap{\raisebox{-0.6ex}{\rule{#1}{0.4pt}}}\ifloesung{\lsgfarbe\small\,#2}\fi}}
\newcommand{\luecke}[1][45mm]{\lsg[#1]{}}
% Lücke bis Zeilenende
\newcommand{\lsgrest}[1]{\leavevmode\unskip\hspace{1.5mm}\ifloesung\rlap{\raisebox{0.8pt}{\lsgfarbe\,#1}}\fi\leaders\hbox{\rule[-0.6ex]{0.5pt}{0.4pt}}\hfill\null\par}
\newcommand{\restzeile}{\lsgrest{}}
% nur in der Lösung sichtbar (Platz bleibt frei), z. B. in Tabellenzellen
\newcommand{\lsgtext}[1]{\ifloesung\leavevmode{\lsgfarbe\small#1}\fi}
\newcommand{\lsgoder}[2]{\ifloesung\leavevmode#1\else#2\fi}
\newcommand{\lueckentext}{\linespread{1.5}\selectfont\raggedright}
\newcommand{\schreibzeile}{\rule[-2.8mm]{0pt}{8mm}}
% Schreiblinien: \linien{n} / \lsglinien{n}{Text}; erste Linie eine volle Zeilenhöhe unter der Frage
\newcommand{\lsglinien}[2]{\par\vspace{0.5mm}\noindent\begin{tikzpicture}
  \path (0,0) (\linewidth-0.4pt,-#1\zeilenabstand-1.5mm);
  \foreach \i in {1,...,#1}{\draw[schreiblinie] (0.2pt,-\i\zeilenabstand) -- (\linewidth-0.2pt,-\i\zeilenabstand);}
  \ifloesung\node[overlay,anchor=north west,inner sep=0pt] at (1mm,-\zeilenabstand+1pt+5mm) {\parbox[t]{\linewidth-2mm}{\lsgfarbe\raggedright\setlength{\baselineskip}{\zeilenabstand}\rule{0pt}{5mm}#2}};\fi
  \end{tikzpicture}\par}
\newcommand{\linien}[1]{\lsglinien{#1}{}}

\usepackage{titlesec}
\newcommand{\website}{https://edu-mrh.de/englisch/w-seminar-time/research-question.html}

% ---------- Handbuch statt Arbeitsblatt: Kopf ohne Name/Datum, Abschnitte mit Linie ----------
\fancyhead[R]{\footnotesize\headlight\textcolor{muted}{Research question \& thesis guide}}
\setcounter{secnumdepth}{0}
\titleformat{\section}{\headfont\large\color{primdark}}{}{0em}{}[\vspace{-1.1ex}\textcolor{prim}{\rule{\linewidth}{0.7pt}}]
\titlespacing*{\section}{0pt}{1.3ex}{-0.9ex}
\titleformat{\subsection}{\headfont\normalsize\color{prim}}{}{0em}{}
\titlespacing*{\subsection}{0pt}{1ex}{-0.2ex}
\setlength{\parskip}{0.28em}
\linespread{0.98}
\newcommand{\stufe}[1]{{\headlight\small\textcolor{muted}{#1}}\par}
\newcommand{\zu}{\,{\normalfont$\rightarrow$}\,}
% Hinweiskasten mit linker Linie (wie in der Vorlage)
\newtcolorbox{hinweis}[1][]{enhanced,colback=primtint,colframe=prim,boxrule=0pt,leftrule=2.5pt,arc=1pt,
  left=2.5mm,right=2.5mm,top=1.2mm,bottom=1.2mm,before skip=1.2mm,after skip=1.2mm,#1}
\newtcolorbox{formel}[1][]{enhanced,colback=warntint,colframe=warn,boxrule=0pt,leftrule=2.5pt,arc=1pt,
  left=2.5mm,right=2.5mm,top=1.2mm,bottom=1.2mm,before skip=1.2mm,after skip=1.2mm,#1}

\begin{document}

\wstitel{Guide · W-Seminar “Time”}{From Topic to Thesis: A W-Seminar Research Guide}%
{Theme: \emph{Time, in all its facets} \,|\, Developing your research question, working thesis and thesis statement}

\begin{minipage}[t]{0.84\linewidth}\vspace{0pt}
A successful W-Seminar paper (\emph{Seminararbeit}) follows a clear intellectual arc: your \textbf{research question} is the compass that guides your investigation, and your \textbf{thesis statement} is the evidence-based answer you reach after exploring the literature – and, if you collect them, your own data. In between, your claim grows: first as a \textbf{working thesis}, then as a \textbf{testable claim} (W4–W7: box \emph{My question – version n}).
\end{minipage}\hfill
\begin{minipage}[t]{0.13\linewidth}\vspace{0pt}\raggedleft
\qrcode[height=15mm]{\website}
\end{minipage}
\par\vspace{1.5mm}
\centerline{\fragepfad{6}}

% ================================================================== STEP 1
\section{Step 1: Craft the Research Question}
\stufe{the starting point · spark question\zu research question · W2–W4}

A research question narrows the vast theme of \emph{Time} into an academic problem you can realistically investigate in 10–15 pages. Start with curiosity, not a fixed opinion – your first \emph{spark question} may be rough.

\begin{enumerate}[leftmargin=1.4em, itemsep=0.05em]
\item \textbf{Choose a facet of “time”:} one disciplinary angle – physics, biology, psychology, history, literature, language …
\item \textbf{Narrow the scope:} restrict your topic by \emph{population, historical era, medium, material or specific variables}.
\item \textbf{Use open-ended interrogatives:} \emph{How}, \emph{Why}, \emph{To what extent}, \emph{In what ways}. Avoid simple yes/no questions.
\item \textbf{Ask one question at a time:} if your question contains an “and”, check whether it hides two.
\end{enumerate}

\begin{hinweis}
\textbf{Weak (yes/no):} Does social media change how teenagers experience time?\\[2pt]
\textbf{Better (open – but still two questions, and very wide):} To what extent does short-form video consumption alter time perception and attention span among adolescents?\\[2pt]
\textbf{Strong (open, focused, researchable):} \emph{To what extent does watching short videos for ten minutes change how long Year 10 students think a following task took?}
\end{hinweis}

\subsection{Check your question: the six markers – or FINER}
\begin{minipage}[t]{0.56\linewidth}\vspace{0pt}
Before committing to a question, check it against the six markers you know from W4:
\par\vspace{0.3ex}
{\small\renewcommand{\arraystretch}{1.08}%
\begin{tabularx}{\linewidth}{@{}>{\bfseries}l X@{}}
Focused & Who or what exactly, where, when?\\
Open & No yes/no, no either/or.\\
Researchable & Do sources exist that you can get and read?\\
Testable & Could a result answer it – either way?\\
Feasible & 10–15 pages, by the deadline, with your means?\\
Relevant & Does someone want to know the answer?\\
\end{tabularx}}
\end{minipage}\hfill
\begin{minipage}[t]{0.41\linewidth}\vspace{0pt}
\begin{hinweis}[colback=sectint,colframe=sec]
{\small\textbf{The same idea, another name: FINER.} Research guides often use this checklist:\par\vspace{0.3ex}
\textbf{F}easible – sources, time\hfill\textcolor{muted}{→ feasible, researchable}\\
\textbf{I}nteresting – to you and others\hfill\textcolor{muted}{→ relevant}\\
\textbf{N}ovel – a fresh angle\hfill\textcolor{muted}{→ relevant}\\
\textbf{E}thical – consent, permission\hfill\textcolor{muted}{→ feasible}\\
\textbf{R}elevant – connected to \emph{Time}\hfill\textcolor{muted}{→ relevant}\par\vspace{0.3ex}
FINER takes for granted what our markers spell out: \emph{focused}, \emph{open}, \emph{testable}.}
\end{hinweis}
\end{minipage}

% ================================================================== STEP 2
\section{Step 2: Conduct Preliminary Research}
\stufe{the bridge · W5}

Before formulating a thesis, survey existing scholarship, historical sources or scientific studies.
\begin{itemize}[leftmargin=1.4em, itemsep=0.05em]
\item \textbf{Map the debate:} identify key concepts, competing theories and unresolved questions in your facet of time – and which methods others used.
\item \textbf{Ask your refinery questions} (W5): what must be defined, explained, put into context, tested and evaluated so that your thesis can stand?
\item \textbf{Follow the evidence:} let credible sources shape your claim rather than cherry-picking quotes to fit a preconceived opinion. Note what speaks \emph{against} your idea, too.
\item \textbf{Few sources?} Look at neighbouring fields, the primary material itself – or ask an expert (M3).
\end{itemize}

% ================================================================== STEP 3
\section{Step 3: Write the Working Thesis}
\stufe{your first answer · research question\zu working thesis · W4–W5}

Your thesis statement is a concise, one-to-two-sentence central claim that directly answers your research question; it acts as an argumentative roadmap for your reader. As long as your research is still running, it is your \textbf{working thesis}: your best answer so far. A simple start: write “I expect that … because …” – then delete “I expect that”.

\begin{formel}
\textbf{Reliable formula:} \quad \textbf{[specific subject]} + \textbf{[debatable claim]} + \textbf{[because of key supporting reasons]}\par\vspace{0.4ex}
{\small The reasons must be something your paper actually examines. \emph{“… because it aligns class schedules with delayed melatonin release”} – you cannot measure melatonin; \emph{“… because students sleep longer”} – that you can ask. The deeper mechanism belongs in your discussion, backed by sources.}
\end{formel}

A good working thesis \textbf{answers} the question directly, is \textbf{arguable} (not a fact, not a taste), \textbf{specific} (group, material, direction), \textbf{testable} (a result could prove it wrong) and \textbf{short}.

% ================================================================== STEP 4
\section{Step 4: Make It Testable}
\stufe{the claim meets the evidence · working thesis\zu testable claim · W6–W7}

A thesis is only as good as the evidence you can actually get – match your claim to it. With your own data (survey, test, experiment) it becomes a \emph{hypothesis} with variables. Always note which result would prove you wrong.

\subsection{Examples across facets of “time”}
{\small\renewcommand{\arraystretch}{1.12}%
\begin{tabularx}{\linewidth}{@{}>{\raggedright\arraybackslash}p{24mm} >{\raggedright\arraybackslash}X >{\raggedright\arraybackslash}X >{\raggedright\arraybackslash}p{34mm}@{}}
\arrayrulecolor{rule}\toprule
\textbf{Facet} & \textbf{Research question} & \textbf{Working thesis} & \textbf{Tested with}\\\midrule
{\headfont\color{primdark}Chronobiology} & How does a later school start affect how long Year 11 students sleep on school nights? & Starting school one hour later lengthens school-night sleep, because adolescents’ sleep phase shifts later. & own data: survey at two schools with different start times\\\midrule
{\headfont\color{primdark}History} & How did local communities in 19th-century Germany respond to standardised railway time? & Resistance to railway time was not merely technical conservatism: it replaced local and religious timekeeping with industrial and state control. & historical sources: newspapers, petitions, council records\\\midrule
{\headfont\color{primdark}Literature} & How does non-linear narrative structure in modernist fiction reflect subjective memory? & Non-linear structures mirror memory by ordering scenes by emotional association rather than chronology – psychological time outweighs clock time. & close reading and coding of scene order in two novels\\\midrule
{\headfont\color{primdark}Psychology} & To what extent does the feeling that time speeds up with age depend on how it is measured? & The age effect appears when people look back over years, not when they judge the present. & studies as data: 14 published studies, grouped by measure\\\bottomrule
\end{tabularx}}

\subsection{Choose the option that fits you}
Of your one to three options (W4), prefer the one where you can say \emph{yes} to most of these: I will still care about it in ten months · I can get and read the sources · the evidence is within reach – own data, studies or texts · it suits my skills (statistics, close reading, languages, Python …) · I know which result would prove it wrong · it fits 10–15 pages. If one is a clear \emph{no}, fix that point first – or drop the option.

% ================================================================== STEP 5
\section{Key Rules for a Strong W-Seminar Paper}
\stufe{testable claim\zu thesis statement · the paper}

\begin{itemize}[leftmargin=1.4em, itemsep=0.05em]
\item \textbf{It must be arguable:} a thesis cannot be an undisputed fact (\emph{“Einstein developed the theory of relativity”}). Another researcher must be able to examine the same evidence and debate your interpretation.
\item \textbf{It must be specific:} replace vague claims like \emph{“Time plays an important role in society”} with exact mechanisms, texts or variables.
\item \textbf{It must be testable:} a claim no result could contradict is a belief; a wish (\emph{“Schools should start later”}) needs a claim behind it.
\item \textbf{Expect evolution:} treat your claim as a working thesis and refine it version by version (W4–W7: \emph{what changed – and why?}) – a hallmark of strong academic work. The \textbf{thesis statement} comes last, after your results: it opens your introduction and answers your question in the conclusion. If the data do not support your working thesis, that is a result, too – report it and discuss why.
\end{itemize}

\subsection{Example: one question, six versions \textcolor{muted}{(as in the worksheet solutions)}}
{\small\renewcommand{\arraystretch}{1.02}%
\begin{tabularx}{\linewidth}{@{}>{\raggedright\arraybackslash}p{19mm} >{\raggedright\arraybackslash}X >{\raggedright\arraybackslash}p{48mm}@{}}
\arrayrulecolor{rule}\toprule
\textbf{version} & \textbf{question / claim} & \textbf{what changed – and why}\\\midrule
0 · W2 spark & Why do waits feel so long? & curiosity, no focus yet\\
1 · W4 & Does music change how long a wait feels? – \emph{Music makes waits feel shorter.} & a focus – but still yes/no\\
2 · W5 & How does the tempo of background music change how long a wait feels? – \emph{With fast music, a 3-minute wait feels shorter than with slow music.} & sources: the tempo matters, not music as such\\
3 · W6 & \emph{Fast music (140 bpm) leads to lower estimates of a 3-minute wait, in minutes, than slow music (60 bpm).} & made measurable: fixed wait, two tempos, a unit\\
4 · W7 & same claim; two clearly different tempos, about 64 students per group & pilot: tempos felt too similar, too few people\\
the paper & \emph{For teenagers, fast background music makes a 3-minute wait feel almost half a minute shorter than slow music.} & thesis statement, backed by the results (d = 0.73)\\\bottomrule
\end{tabularx}}

\end{document}
